% ch-eu-proxies -- the opaque regime: proxies from deal marks and the
% selection problem (imported from the handbook's US chapter; body of
% "When Fundamentals Are Latent: Proxies and Selection")

This chapter is the hinge of the transparent-regime part. Its place on the decision map of
\cref{ch01:sec:map} is to explain why the rest of the document exists. The two preceding
chapters assumed fundamentals arrive on a statutory clock: every firm, every fiscal year, a
filed income statement. In the US-centric PitchBook universe of \cref{ch:data} that clock does
not exist. The vendor delivers instead the firm's \emph{transaction record}: a sequence of deal
events, each carrying marks --- a valuation, a round classification, an employee count,
occasionally a revenue or EBITDA figure. Between events we learn nothing. This chapter
re-derives the transparent-regime pipeline under that observation model. Three things change:
the \emph{target} (EBITDA growth must be replaced by proxies built from marks), the \emph{unit
of analysis} (the firm-year panel row is replaced by the consecutive-deal pair), and --- most
consequentially --- the \emph{bias structure}: outcomes are observed only when the firm
transacts again, so every naive estimate is filtered through survival-to-next-round. Against
the defect ledger of \cref{tab:ch02:defects}, this is where (C5) and (C3) bind hardest: the
covariates are the LOCF process of \cref{eq:ch02:locf} by construction, and the outcomes
inherit the immortal-firm selection. The repair for that selection is a model of the
deal-arrival process itself, which is why the continuous-time machinery of the next part
completes this chapter.

\section{The observation model}\label{chpx:sec:obs}

Let $S_{i,t} \in \R^d$ denote the latent state of firm $i$: fundamentals (revenue, margin
trajectory, unit economics) plus soft state (team quality, runway). In the transparent regime we
observed a noisy annual census of $S$. Here we observe marks only at the firm's own deal times
$T_{i,1} < T_{i,2} < \cdots$:
\begin{equation}\label{eq:us-marks}
m_{i,k} \;=\; \big(V_{i,k},\ \code{round\_type}_{i,k},\ E_{i,k},\
R_{i,k},\ B_{i,k}\big),
\end{equation}
where $V$ is post-money valuation (populated only for priced rounds), $E$ the reported employee
count, and $R$, $B$ the revenue and EBITDA fields --- populated for roughly $45\%$ and $25\%$ of
deal rows respectively in typical vendor coverage (\cref{ch:data},
\cref{app:columns}). The marks are noisy, truncated views of $S_{i,T_{i,k}}$ in the classical
errors-in-variables sense \citep{carroll2006}: valuation confounds fundamentals with the market
multiple, employees confound productivity with cash deployment, and the revenue field is
populated selectively.

What the modeler works with between deals is exactly the LOCF panel formalized once in
\cref{ch02:sec:formal}: the carried-forward mark vector $\tilde S_{i,t}$ of \cref{eq:ch02:locf}
together with the staleness clock $\staleage_i(t) = t - \tau_i(t)$. We do not restate that
formalization here; \cref{ch:stale} owns it, including the proof that the pair $(\tilde S_i,
\staleage_i)$ is predictable and therefore a legitimate covariate, and the analysis of what
staleness does to fitted effects. Two consequences matter for this chapter. The staleness clock
carries real information: a valuation carried forward $4$ years is a different object from one
carried forward $4$ months, and models must be allowed to discount stale marks. And the deeper
problem is \emph{when} the panel refreshes:

\begin{warning}[Sampling times are informative]\label{warn:us-mnar}
The observation times $T_{i,k}$ are not an exogenous measurement schedule; they are outcomes of
the firm's own trajectory. A firm raises when it is doing well (an up round at a rich multiple)
or when it must (runway exhaustion, a down round); firms that are neither raise nothing and go
dark. Formally, the event indicator $\dd N_i(t)$ is not independent of $S_{i,t}$ given the past:
the missingness mechanism of the LOCF panel is \emph{missing-not-at-random}. Any panel method
that treats $\tilde{S}_{i,t}$ as if generated by random inspection times --- imputation, mixed
models, naive regressions on the LOCF columns --- inherits a bias whose sign and size depend on
the unmodelled sampling mechanism. The only honest formulations either (i) condition explicitly
on the event frame (\cref{chpx:sec:event}) and correct for selection
(\cref{chpx:sec:selection}), or (ii) model the sampling process jointly with the state
\citep{rizopoulos2012,tsiatisdavidian2004} --- the escalation rung of \cref{ch:stale}, built on
the intensity machinery of \cref{ch:poisson,ch:hawkes}.
\end{warning}

\section{What replaces EBITDA growth}\label{chpx:sec:proxies}

The transparent-regime success labels of \cref{ch:eufund} are defined on filed fundamentals. We
now construct their surrogates, all defined on a pair of consecutive deals $(T_{i,k},
T_{i,k+1})$ with gap $\Delta_{i,k} = T_{i,k+1} - T_{i,k}$.

\paragraph{(a) Implied valuation growth.} When both endpoints are priced rounds,
\begin{equation}\label{eq:us-gval}
g^{\mathrm{val}}_{i,k} \;=\; \frac{\log V_{i,k+1} - \log
V_{i,k}}{\Delta_{i,k}},
\end{equation}
an annualized log growth rate. Writing valuation as fundamentals times a sector multiple, $\log
V_{i,k} = \log F_{i,T_{i,k}} + \log m_{s(i),T_{i,k}} + \varepsilon_{i,k}$, gives the
decomposition
\begin{equation}\label{eq:us-decomp}
g^{\mathrm{val}}_{i,k}
\;=\; \underbrace{g^{F}_{i,k}}_{\text{fundamentals growth}}
\;+\; \underbrace{g^{m}_{i,k}}_{\text{multiple expansion}}
\;+\; \underbrace{\nu_{i,k}}_{\text{pricing noise}} .
\end{equation}
Only $g^{F}$ is a success signal; $g^{m}$ is a market-wide repricing that the firm did not earn.
We identify the multiple component by the public-market-comparables logic: regress
$g^{\mathrm{val}}$ on sector $\times$ time-bucket fixed effects plus macro controls (rate
changes, public-index sector returns over the gap); the fitted cell means estimate $g^{m}$, and
the residual is the firm-level proxy $\hat{g}^{F}_{i,k}$. The fixed effects must be estimated on
all pairs in the cell, not per firm, or they absorb the signal.

\paragraph{(b) Round direction as an ordinal outcome.} The field
\code{vc\_round\_up\_down\_flat} classifies each round against the previous one as
$\{\text{down}, \text{flat}, \text{up}\}$. It is coarse but far better populated than paired
valuations (it is defined even when one endpoint is unpriced) and robust to level noise in $V$.
We treat it as an ordinal target (proportional-odds or ordinal gradient boosting); its drawback
is the vendor's flat band, which is a definitional artifact rather than an economic category.

\paragraph{(c) Employee growth.} $g^{\mathrm{emp}}_{i,k} = (\log E_{i,k+1} - \log
E_{i,k})/\Delta_{i,k}$ is a real-activity proxy available for most pairs. Its bias is
mechanical: headcount growth immediately after a raise reflects cash deployment, not
demonstrated productivity, so it leads fundamentals optimistically by construction.

\paragraph{(d) Reported revenue growth.} \code{revenue\_growth\_since\_last\_deal} is the most
direct surrogate when populated, but with $\sim\!45\%$ endpoint coverage the pair-level coverage
is roughly $20\%$, and reporters skew to larger, later-stage firms --- the missingness is again
not at random.

\begin{table}[htbp]
\centering
\small
\begin{tabular}{p{3.1cm}p{2.1cm}p{2.5cm}p{4.2cm}}
\toprule
Proxy & Pair coverage & Noise & Dominant bias \\
\midrule
$g^{\mathrm{val}}$ \cref{eq:us-gval} & moderate (both rounds priced) &
high (negotiation, multiple) & upward; confounds multiple expansion
unless decomposed via \cref{eq:us-decomp} \\
\code{vc\_round\_}\newline\code{up\_down\_flat} & high & low but coarse
& flat band definitional; censors magnitude \\
$g^{\mathrm{emp}}$ & high & moderate & upward after raises (cash-fueled
hiring) \\
\code{revenue\_growth\_}\newline\code{since\_last\_deal} & low
($\sim\!20\%$) & low given reported & reporters skew large/late-stage;
MNAR \\
\bottomrule
\end{tabular}
\caption{Growth proxies on the consecutive-deal pair frame. All four
are additionally subject to the survival-to-next-round selection of
\cref{chpx:sec:selection}, which no choice among them avoids.}
\label{chpx:tab:proxies}
\end{table}

In practice we fit all four and read them jointly: agreement across proxies with different bias
directions is the closest available substitute for a filed income statement.

\section{Supervised learning on the event frame}\label{chpx:sec:event}

The prediction unit is the ordered pair $(i,k) \to (i,k+1)$: one row per consecutive-deal pair,
outcome one of the proxies of \cref{chpx:sec:proxies}. Features come in two blocks with
different epistemic status.

\emph{Firm-history block}, measurable at $T_{i,k}$: the full mark history through the $k$-th
deal (valuation path, round sequence, employee path, deal sizes, investor counts), the staleness
pattern, firm age, sector, and gap lengths between \emph{past} deals. Nothing dated inside
$(T_{i,k}, T_{i,k+1}]$ from the firm's own record enters --- there is nothing there to use
except the second deal itself, which is the outcome's timestamp.

\emph{Gap macro block}, aggregates of exogenous series over the gap: mean policy rate, change in
CPI year-on-year, change in unemployment, and sector deal heat (sector deal counts per week,
normalized by trailing average), all built from point-in-time vintages per \cref{prot:ch02:pit}.
Because these are external to the single firm, conditioning on them does not leak the firm's
outcome; but it does change the estimand. With the gap block included, the regression is an
\emph{attribution} model --- growth conditional on the realized macro path and the realized gap
length. For deployment as a forecast at $T_{i,k}$, the gap block must be replaced by scenario
paths or integrated over a gap-length model; the gap length $\Delta_{i,k}$ is itself informative
(short gaps select hot firms), and we return to precisely this point below.

The learner is the \cref{ch:eufund} stack unchanged: gradient boosting with blocked
cross-validation (folds split by firm \emph{and} time, under the constitution of
\cref{ch07:sec:temporal}), early stopping on the temporal holdout, and conformalized quantile
intervals for uncertainty; we do not repeat that machinery here. Two adaptations matter. First,
conformal calibration sets must be blocked by firm, since pairs from one firm share latent
quality and violate exchangeability across naive random splits. Second, once the selection
weights of the next section are introduced, calibration must be weight-aware (weighted conformal
quantiles), otherwise coverage holds for the selected population only.

\section{The selection problem}\label{chpx:sec:selection}

This section is the core of the chapter. A pair $(i,k)$ exists in the training frame only if the
firm raised again within the observation window: define
\begin{equation}\label{eq:us-selind}
D_{i,k} \;=\; \ind{T_{i,k+1} \le \tau_{\mathrm{end}}},
\end{equation}
where $\tau_{\mathrm{end}}$ is the data cutoff. Every proxy of \cref{chpx:sec:proxies} is
observed conditional on $D_{i,k} = 1$: we see the growth of firms \emph{that came back}. The
observed distribution is $P(g \st D{=}1)$, not $P(g)$, and the gap between them has an exact
expression:
\begin{equation}\label{eq:us-selbias}
\E[g \st D{=}1] - \E[g]
\;=\; \frac{\Cov\!\big(g,\ \ind{D{=}1}\big)}{\Prob(D{=}1)} .
\end{equation}
Both selection channels have positive covariance with \emph{appearing} in the frame but opposite
implications for the observed growth mix: winners raise again (pulling observed $g$ up), and
cash-constrained firms raise again at down rounds (pulling it down), while the quiet failures
--- firms that never transact again --- contribute nothing at all. The left tail of $P(g)$ is
therefore systematically missing: shutdown is the one outcome that generates no second deal.
Estimates of mean growth, of feature effects, and of ranking quality are all computed on the
surviving subpopulation --- and reading the survivors' serial-raising pattern as firm-level
state dependence rather than selection is the classical trap \citep{heckmanborjas1980}.

\paragraph{Bounds.} If $g$ is bounded to $[g_{\mathrm{lo}}, g_{\mathrm{hi}}]$ (after
winsorization, or by construction for the ordinal outcome), then with $p = \Prob(D{=}1)$,
\begin{equation}\label{eq:us-bounds}
\E[g] \;\in\; \big[\, p\,\E[g \st D{=}1] + (1-p)\,g_{\mathrm{lo}},\;\;
p\,\E[g \st D{=}1] + (1-p)\,g_{\mathrm{hi}} \,\big],
\end{equation}
the worst-case (no-assumption) bounds. With $p$ well below one --- deal records concentrate on
a minority of serially raising firms (\cref{chpx:fig:deals-per-firm}) --- these bounds are
wide. They are still worth reporting: they state what the data alone say, and any point
estimate inside them relies on an assumption.

\begin{assumption}[Selection on observed history]\label{ass:us-soh}
$g_{i,k} \indep D_{i,k} \st \Hist_{T_{i,k}}$, where $\Hist_{T_{i,k}}$ collects the firm's mark
history through deal $k$ and the macro state: given everything observed up to the first deal of
the pair, whether the firm raises again carries no further information about its growth.
\end{assumption}

\paragraph{The practical fix: inverse-probability-of-observation weights.} Under
\cref{ass:us-soh}, define the observation probability
\begin{equation}\label{eq:us-pi}
\pi_{i,k} \;=\; \Prob\big(D_{i,k}{=}1 \st \Hist_{T_{i,k}}\big)
\;=\; 1 -
\exp\!\Big(-\!\int_{T_{i,k}}^{\tau_{\mathrm{end}}}
\lambda_i(u)\dd u\Big),
\end{equation}
where $\lambda_i$ is the firm's next-deal intensity given its history. Weighting each observed
pair by $1/\hat{\pi}_{i,k}$ (stabilized and trimmed in the usual way) makes the weighted frame
representative of all at-risk firms: $\E[D g / \pi] = \E[g]$. This is inverse-probability
weighting for informative censoring in the sense of \citet{hernan2020}, with the hazard-model
lineage of \citet{cox1972}; validity requires $\pi_{i,k} > 0$ and a well-specified intensity,
and residual selection on latent quality beyond $\Hist$ is not repaired --- which is why we
report \cref{eq:us-bounds} alongside. The selection identities \cref{eq:us-selbias,eq:us-bounds}
and the unbiasedness of the weighted mean are proved in \cref{app:proofs}.

\begin{remark}[The bridge to the machinery]\label{rem:us-bridge}
Equation \cref{eq:us-pi} is the structural link between this document's two halves. The
quantity $\lambda_i(u)$ --- the conditional intensity of the firm's deal-arrival process --- is
exactly what the next part calibrates: inhomogeneous Poisson with macro covariates
(\cref{ch:poisson}), self-exciting and inhibiting extensions (\cref{ch:hawkes}), and the
censored variants that untrusted dates force (\cref{ch:censoring}). The count model is
therefore not a separate deliverable that happens to share the data: it \emph{repairs the
growth regression}. A team that skips the point-process layer does not merely lose the
deal-timing forecast; it loses the only principled weights for its growth model. The combined
pipeline --- intensity fit, weights, weighted growth regression --- is to be validated on the
synthetic world of \cref{ch02:sec:synthetic}, where $P(g)$ is known, with results reported in
\cref{ch:results-synth}.
\end{remark}

\begin{figure}[htbp]
\centering
\includegraphics[width=.85\textwidth]{fig_deals_per_firm}
\caption{Rank-size plot of deals per firm in the two synthetic regimes
(illustrative; the generator reproduces the concentration typical of
deal records). A small minority of serially raising firms contributes
most consecutive-deal pairs, while the modal firm contributes one deal
and zero pairs. The pair frame therefore oversamples the firms with the
highest historical deal hazard; unweighted estimates describe serial
raisers, not the tracked universe. Firm-level cluster bootstrap is
mandatory for all standard errors on this frame.}
\label{chpx:fig:deals-per-firm}
\end{figure}

\section{Macro factors as drivers}\label{chpx:sec:macro}

With firm-level outcomes this sparse and selected, much of the macro question --- do rate
changes move private-market activity and pricing? --- is better asked at the sector level. The
workhorse is a sector-exposure (shift-share style) design. Let $y_{s,w}$ be a sector-week
outcome (log deal count, or the cell-mean valuation growth from \cref{eq:us-decomp}), $\Delta
r_w$ the rate change, and $e_s$ a \emph{predetermined} exposure measure (share of pre-revenue
firms in the sector, a duration proxy for how far cash flows sit in the future). We fit
\begin{equation}\label{eq:us-shiftshare}
y_{s,w} \;=\; \alpha_s + \gamma_w + \delta\,(e_s \cdot \Delta r_w) +
u_{s,w},
\end{equation}
with sector and time fixed effects and two-way clustered standard errors. The time effect
$\gamma_w$ absorbs every common shock --- including the \emph{average} effect of the rate change
itself --- so $\delta$ is identified purely from differential sensitivity: do high-exposure
sectors move more when rates move?

Three interpretation limits apply. First, the common component of the macro effect is not
identified by design; \cref{eq:us-shiftshare} estimates relative sensitivities only. Second,
latent risk appetite co-moves with rates; if its sector loadings correlate with $e_s$, the
differential estimate is contaminated too. Third, exposure must be predetermined and not itself
an outcome of recent sector performance. A design-based reading in the spirit of
\citet{imbens2015} is available only under an explicit exogeneity assumption on the exposure
shares, which we state and plan to stress-test on the synthetic testbed
(\cref{ch:results-synth}) rather than assume silently. Where valuation marks are the outcome,
all of \cref{chpx:sec:selection}'s caveats apply on top.

The forward pointer of the section: deal \emph{counts} are complete (a deal either appears in
the feed or does not), while marks are sparse and selection-laden. The cleanest macro inference
therefore lives at the count layer --- intensities modulated by macro covariates, the
covariate-NHPP baseline that \cref{ch:poisson} builds and that every excitation story must
beat.

\section{Decision table and the bridge out}\label{chpx:sec:decision}

\begin{table}[htbp]
\centering
\small
\begin{tabular}{p{2.5cm}p{4.6cm}p{5.6cm}}
\toprule
Aspect & Transparent track (\cref{ch:eufund,ch:eucausal}) & Latent
adaptation (this chapter) \\
\midrule
Target & EBITDA/revenue growth from filings & proxy set: $\hat{g}^{F}$
from \cref{eq:us-decomp}, ordinal round direction, employee growth,
reported revenue growth \\
Unit & firm $\times$ fiscal year & consecutive-deal pair $(i,k)$ \\
Feature clock & fiscal-year as-of with publication lags & event-time
as-of with staleness clock (\cref{ch02:sec:formal}, \cref{ch:stale}) \\
Dominant bias & leakage via publication timing & survival-to-next-round
selection; MNAR sampling times \\
Primary fix & point-in-time joins, blocked CV & IP-of-observation
weights from the next-deal intensity \cref{eq:us-pi}; worst-case bounds
\cref{eq:us-bounds} \\
Uncertainty & conformal intervals, firm-blocked & weight-aware
conformal, firm-cluster bootstrap \\
Macro layer & covariates in the growth model & shift-share
differentials \cref{eq:us-shiftshare}; count-layer inference in
\cref{ch:poisson} \\
\bottomrule
\end{tabular}
\caption{Method mapping between observation regimes. Every row of the
right column degrades gracefully to the left column when fundamentals
become observable; the converse direction is the content of this
chapter.}
\label{chpx:tab:decision}
\end{table}

\Cref{chpx:tab:decision} summarizes the chapter and closes the transparent-regime part with an
explicit hand-off. The selection problem exposed here is not patched within this chapter's
toolkit: its repair \emph{is} the hazard and point-process machinery that follows --- the
next-deal intensity of \cref{eq:us-pi} is fitted by the Poisson and Hawkes layers of
\cref{ch:poisson,ch:hawkes}, and the discrete-time face of the same repair (censored buckets
labeled, never dropped; eligibility masks) is the bucketed-panel discipline of
\cref{ch:buckets}. Likewise, the staleness of carried-forward marks is not formalized here ---
the LOCF observation model belongs to \cref{ch02:sec:formal}, and its consequences
(attenuation, informative observation times, the remedy ladder) are the subject of
\cref{ch:stale}. Nothing in this chapter's adaptation is exotic: each entry is a
transparent-regime method re-derived under a harsher observation model, with one genuinely new
ingredient --- the deal-arrival intensity --- that the next part constructs.

\begin{decision}{D-EU.3 --- no unweighted pair-frame estimate ships}
Proxy outcomes on the consecutive-deal pair frame are always reported as a joint read across the
four proxies of \cref{chpx:tab:proxies}, and any population-level statement (mean growth,
feature effect, ranking quality) is computed with inverse-probability-of-observation weights
\cref{eq:us-pi} derived from a fitted next-deal intensity, flanked by the worst-case bounds
\cref{eq:us-bounds}; unweighted estimates are labeled as describing the serial-raiser
subpopulation and never leave the diagnostic tier. Rationale: \cref{eq:us-selbias} is an
identity, not a risk --- the selected and full-population means differ whenever growth covaries
with returning to market, and the frame concentration of \cref{chpx:fig:deals-per-firm} makes
the covariance large. Reversal trigger: if fitted observation probabilities are uniformly
near one within a segment (e.g.\ late-stage, short-horizon pairs), the weights degrade to the
identity there and this box's requirement reduces to reporting that fact.
\end{decision}

\begin{codehooks}
The pair-frame builder is a planned module \modrepo{pairs/}: it consumes the event list and LOCF
features emitted by the loaders (\code{synthetic/loaders.py} now, \modrepo{data/pitchbook.py}
later), enforces the firm-history/gap-macro block separation of \cref{chpx:sec:event}, and emits
$(i, k, \Delta_{i,k}, g^{\cdot}, D_{i,k})$ rows with staleness ages attached. The observation
weights \cref{eq:us-pi} are computed from the intensity fitters of the next part ---
\modrepo{eventtime/} for firm-level gap models, \modrepo{sector\_hazard.py} for the sector layer
--- so the module depends on \cref{ch:poisson,ch:hawkes} deliverables by design. The
multiple-decomposition regression \cref{eq:us-decomp} and the shift-share design
\cref{eq:us-shiftshare} share the sector-week frame of \modrepo{sector\_stability.py}.
Validation of the combined weights-plus-regression pipeline against known $P(g)$ will run on
\code{synthetic/} ground truth (regimes A and B), with results reported in
\cref{ch:results-synth}.
\end{codehooks}
